\documentclass[12pt,a4paper]{article}

% ══════════════════════════════════════════════════════════════
% Preâmbulo compartilhado pelos documentos da P2
% Reaproveita o estilo de docs/documento-unificado.tex
% ══════════════════════════════════════════════════════════════

\usepackage[utf8]{inputenc}
\usepackage[T1]{fontenc}
\usepackage[brazil]{babel}

\usepackage[top=2.5cm, bottom=2.5cm, left=3cm, right=2cm]{geometry}
\usepackage{setspace}
\onehalfspacing
\setlength{\parindent}{0pt}
\setlength{\parskip}{6pt}

\usepackage{lmodern}
\usepackage{microtype}
\usepackage{amssymb}
\usepackage{amsmath}
\usepackage{pifont}
\newcommand{\cmark}{\ding{51}}
\newcommand{\xmark}{\ding{55}}

\usepackage[table,dvipsnames]{xcolor}
\definecolor{accent}{HTML}{2C7A4B}
\definecolor{accentdim}{HTML}{1E5934}
\definecolor{accentlight}{HTML}{E8F4ED}
\definecolor{tabheader}{HTML}{2C3E50}
\definecolor{tabrow}{HTML}{F5F6FA}
\definecolor{linkblue}{HTML}{2980B9}
\definecolor{codebg}{HTML}{F5F5F5}
\definecolor{codeborder}{HTML}{D0D0D0}
\definecolor{ndvilow}{HTML}{8B4513}
\definecolor{ndvimid}{HTML}{DAA520}
\definecolor{ndvihigh}{HTML}{006400}
\definecolor{warnyellow}{HTML}{F39C12}
\definecolor{successgreen}{HTML}{27AE60}
\definecolor{errorred}{HTML}{E74C3C}

\usepackage{booktabs}
\usepackage{tabularx}
\usepackage{multirow}
\usepackage{makecell}
\usepackage{float}
\usepackage{graphicx}
\usepackage{caption}
\usepackage{array}

\usepackage{enumitem}
\setlist[itemize]{itemsep=2pt, parsep=0pt, topsep=4pt, leftmargin=*}
\setlist[enumerate]{itemsep=2pt, parsep=0pt, topsep=4pt, leftmargin=*}

\usepackage[colorlinks=true, linkcolor=linkblue, urlcolor=linkblue, citecolor=linkblue]{hyperref}

\usepackage{fancyhdr}
\pagestyle{fancy}
\fancyhf{}
\fancyhead[L]{\small\textcolor{gray}{EasyGis}}
\fancyhead[R]{\small\textcolor{gray}{P2 — Gerência e Manutenção de Software · Prof.\ Mauro · 2026.1}}
\fancyfoot[C]{\thepage}
\renewcommand{\headrulewidth}{0.4pt}

\usepackage{titlesec}
\titleformat{\section}
  {\Large\bfseries\color{accentdim}}
  {\thesection.}{0.5em}{}
\titleformat{\subsection}
  {\large\bfseries\color{accent}}
  {\thesubsection}{0.5em}{}
\titleformat{\subsubsection}
  {\normalsize\bfseries}
  {\thesubsubsection}{0.5em}{}

% Listings (código)
\usepackage{listings}
\lstdefinestyle{shellstyle}{
  basicstyle=\ttfamily\small,
  backgroundcolor=\color{codebg},
  frame=single,
  rulecolor=\color{codeborder},
  framesep=4pt,
  xleftmargin=8pt,
  xrightmargin=4pt,
  breaklines=true,
  showstringspaces=false,
  keepspaces=true,
  literate={─}{{-}}1 {├}{{|}}1 {│}{{|}}1 {└}{{`}}1
}
\lstdefinestyle{pythonstyle}{
  basicstyle=\ttfamily\footnotesize,
  language=Python,
  backgroundcolor=\color{codebg},
  frame=single,
  rulecolor=\color{codeborder},
  framesep=4pt,
  xleftmargin=8pt,
  xrightmargin=4pt,
  breaklines=true,
  keywordstyle=\color{accent}\bfseries,
  stringstyle=\color{ndvimid},
  commentstyle=\color{gray}\itshape,
  showstringspaces=false,
  keepspaces=true
}
\lstset{style=shellstyle}

% Caixas para código inline
\newcommand{\code}[1]{\texttt{\colorbox{codebg}{\strut#1}}}

% Caixas de destaque
\usepackage{tcolorbox}
\tcbuselibrary{skins, breakable}

\newtcolorbox{infobox}[1][]{%
  colback=accentlight, colframe=accent,
  boxrule=0.6pt, arc=2pt,
  left=8pt, right=8pt, top=4pt, bottom=4pt,
  fonttitle=\bfseries, breakable, #1
}

\newtcolorbox{warnbox}[1][]{%
  colback=yellow!8, colframe=warnyellow,
  boxrule=0.6pt, arc=2pt,
  left=8pt, right=8pt, top=4pt, bottom=4pt,
  breakable, #1
}

\newtcolorbox{successbox}[1][]{%
  colback=green!4, colframe=successgreen,
  boxrule=0.6pt, arc=2pt,
  left=8pt, right=8pt, top=4pt, bottom=4pt,
  breakable, #1
}


\title{\textbf{Roteiro da Apresentação P2}\\[6pt]
       \large EasyGis — Distribuição de slides e falas por integrante}
\author{Engenharia de Software · 2026.1}
\date{Maio · 2026}

\begin{document}
\maketitle

\section*{Visão geral}

A apresentação tem \textbf{14 slides} cobrindo o ciclo completo do produto:
problema, solução, escopo, arquitetura, fluxo técnico, demo ao vivo, testes,
resultados, limitações e equipe. Tempo previsto: \textbf{17--18 minutos}.

\vspace{4pt}
A equipe é dividida da seguinte forma:

\begin{center}
\renewcommand{\arraystretch}{1.35}
\begin{tabularx}{\textwidth}{@{}l l X@{}}
\toprule
\textbf{\#} & \textbf{Slide} & \textbf{Apresentador} \\
\midrule
1  & Capa                              & Lucas Benevinuto \\
2  & Problema                          & Lucas Benevinuto \\
3  & Solução                           & Lucas Benevinuto \\
4  & Para quem                         & João Vinícius Castello \\
5  & Backlog (30/35 entregues)         & José Melquíades \\
6  & Arquitetura                       & Vinicius Henrique \\
7  & Stack tecnológica                 & Luis Gustavo \\
8  & Decisões de design (ADRs)         & Luis Gustavo \\
9  & Fluxo do NDVI                     & João Leonardi \\
10 & Demonstração ao vivo              & João Vinícius Castello (condutor) \\
11 & Estratégia de testes              & Lauan Matheus \\
12 & Resultados (28/28, cobertura)     & Lauan Matheus \\
13 & Limitações e próximos passos      & Sammuel Moura \\
14 & Equipe e Perguntas?               & João Leonardi (encerra) + todos \\
\bottomrule
\end{tabularx}
\end{center}

\vspace{6pt}
\textbf{Resumo por integrante:}

\begin{itemize}
  \item \textbf{Lucas Benevinuto} — slides 1, 2, 3 (abertura + problema + solução)
  \item \textbf{João V. Castello} — slides 4, 10 (perfis + condução da demo)
  \item \textbf{José Melquíades} — slide 5 (backlog)
  \item \textbf{Vinicius Henrique} — slide 6 (arquitetura)
  \item \textbf{Luis Gustavo} — slides 7, 8 (stack + ADRs)
  \item \textbf{João Leonardi} — slides 9, 14 (fluxo NDVI + encerramento)
  \item \textbf{Lauan Matheus} — slides 11, 12 (testes + resultados)
  \item \textbf{Sammuel Moura} — slide 13 (limitações + próximos passos)
\end{itemize}

\newpage

\section*{Falas — slide a slide}

\subsection*{Slide 1 — Capa \hfill \textit{Lucas Benevinuto · 30 s}}

\textquotedblleft\,Boa tarde, professor. Somos a equipe que apresenta o
\textbf{EasyGis} — uma plataforma web de agricultura de precisão por
sensoriamento remoto. Nosso trabalho foi entregar uma ferramenta que
transforma imagens Sentinel-2 em decisões sobre uma lavoura, sem custo
para o produtor. Eu sou o Lucas, vou conduzir a abertura.\,\textquotedblright

\subsection*{Slide 2 — Problema \hfill \textit{Lucas Benevinuto · 60 s}}

\textquotedblleft\,Para entender o porquê do projeto, três números:
o Brasil tem \textbf{66 milhões de hectares} plantados, a Agência
Espacial Europeia disponibiliza imagens \textbf{Sentinel-2} de \textbf{10\,m
de resolução a cada 5 dias}, \textbf{de graça}. Apesar disso, ferramentas
prontas que usam esses dados custam \textbf{milhares de reais por ano}
ou exigem que o usuário monte uma stack de GDAL, scripts e SIG no
desktop. O EasyGis ataca exatamente esse gap: aproximar essa imagem
gratuita do produtor.\,\textquotedblright

\subsection*{Slide 3 — Solução \hfill \textit{Lucas Benevinuto · 50 s}}

\textquotedblleft\,A solução é uma \textbf{página web local, sem custo}.
O fluxo é direto: o usuário \textbf{importa um KML} do Google Earth ou
desenha o talhão no mapa; escolhe o \textbf{índice} (NDVI, EVI, SAVI,
NDWI, NDBI ou bandas brutas); e em \textbf{5 a 30 segundos} aparece a
imagem colorida sobreposta ao mapa, com estatísticas e classificação em
zonas de vigor com área em hectares. Tudo isso em uma única tela.
Quem vai detalhar para quem isso serve é o João Castello.\,\textquotedblright

\subsection*{Slide 4 — Para quem \hfill \textit{João V. Castello · 60 s}}

\textquotedblleft\,O EasyGis foi pensado para \textbf{quatro perfis} muito
diferentes, mas que olham para o mesmo mapa. O \textbf{produtor rural}
quer saber onde aplicar amanhã. O \textbf{agrônomo} precisa de um mapa de
variabilidade para prescrever insumos. O \textbf{pesquisador} usa a
ferramenta como sandbox para validar algoritmos contra dados reais. E o
\textbf{estudante} aprende sensoriamento remoto na prática, sem barreira
de licença ou setup. Mesmo produto, leituras diferentes — e isso guiou
nossa priorização. Vou passar para o José falar sobre escopo.\,\textquotedblright

\subsection*{Slide 5 — Backlog \hfill \textit{José Melquíades · 80 s}}

\textquotedblleft\,Sobre o escopo: começamos com \textbf{35 histórias
distribuídas em 6 épicos} — importação, índices, visualização,
classificação, 3D e infraestrutura. Priorizamos por \textbf{MoSCoW}:
\textit{Must Have}, \textit{Should}, \textit{Could} e \textit{Won't
Have}. Entregamos \textbf{28 histórias completas} e \textbf{2 parciais},
o que dá \textbf{85,7\% do escopo planejado}, dentro do prazo do
semestre. Os \textbf{5 \textit{Won't Have}} foram \textbf{conscientes}:
autenticação, persistência em banco, exportação e multi-usuário ficaram
fora porque o foco era validar o \textbf{núcleo de processamento}
geoespacial — e isso foi escolha, não acidente.\,\textquotedblright

\subsection*{Slide 6 — Arquitetura \hfill \textit{Vinicius Henrique · 75 s}}

\textquotedblleft\,Arquitetura: optamos por \textbf{dois processos
separados}. O \textbf{frontend} é um Next.js 16 rodando na porta 3000 —
React 19, TypeScript, Tailwind, Leaflet em 2D e Three.js para o relevo em
3D. O \textbf{backend} é uma API FastAPI em Python 3.12 na porta 8000,
com rasterio, NumPy, SciPy, OpenCV e scikit-learn. A comunicação entre
eles é \textbf{REST com JSON}; as imagens processadas viajam em
\textbf{base64} dentro da resposta. \textbf{Não há banco de dados} — os
KMLs e os produtos Sentinel ficam no filesystem. Essa decisão está
registrada como \textbf{DA-02} no documento de ADRs e é o que torna o
sistema reprodutível: qualquer máquina com Python e Node sobe a
aplicação em minutos.\,\textquotedblright

\subsection*{Slide 7 — Stack tecnológica \hfill \textit{Luis Gustavo · 75 s}}

\textquotedblleft\,Cada peça da stack tem uma razão. \textbf{FastAPI}
porque o Pydantic dá validação de contrato automática e gera a doc
OpenAPI em \texttt{/docs} de graça. \textbf{rasterio} porque é o padrão
de fato sobre o GDAL no ecossistema Python. \textbf{Next.js App Router}
porque combina SSR e API routes no mesmo projeto, sem expor o
filesystem ao cliente. \textbf{Leaflet} em vez de Mapbox porque não há
custo por tile carregado. \textbf{Three.js} para visualização 3D nativa
em WebGL. E \textbf{uv} no Python — gerenciador de pacotes
\textbf{dez vezes mais rápido que o pip}, com lockfile determinístico
para garantir reprodução em qualquer máquina. Nada foi escolhido por
hype — toda decisão está justificada.\,\textquotedblright

\subsection*{Slide 8 — Decisões de design \hfill \textit{Luis Gustavo · 80 s}}

\textquotedblleft\,Documentamos \textbf{11 ADRs} no repositório, e
quero destacar três que são chave para entender por que o sistema é
robusto.

\textbf{DA-06 — Separação processador / rotas HTTP.} A lógica de
cálculo dos índices vive numa classe \texttt{SentinelProcessor} que
\textbf{não conhece HTTP}. Isso permite testar a matemática sem subir
servidor — são 28 testes unitários em menos de 1 segundo.

\textbf{DA-07 — Tipagem fim-a-fim.} O Pydantic no backend é
espelhado por interfaces TypeScript no frontend. Um erro de contrato
vira \textbf{422 ou erro de build}, nunca \textbf{500 silencioso} em
produção.

\textbf{DA-09 — Suavização gaussiana preservando NaN.} Um filtro
gaussiano direto \textbf{sangra valores} para fora do polígono e cria
halo nas bordas. Resolvemos dividindo também pela máscara binária — o
detalhe técnico que separa o resultado correto do
plausível-mas-errado.\,\textquotedblright

\subsection*{Slide 9 — Fluxo do NDVI \hfill \textit{João Leonardi · 75 s}}

\textquotedblleft\,Para mostrar como tudo isso se encadeia, vamos pelo
caminho do NDVI em \textbf{5 etapas}.

\textbf{(01)} O usuário seleciona o talhão e o índice no frontend.
\textbf{(02)} O Next.js dispara \texttt{POST /calculate-index} para o
backend. \textbf{(03)} O \texttt{SentinelProcessor} usa rasterio para
ler as bandas \textbf{B08 (infravermelho próximo)} e \textbf{B04
(vermelho)} do produto \texttt{.SAFE}. \textbf{(04)} O pyproj reprojeta
do WGS84 do KML para o UTM do raster. \textbf{(05)} NumPy faz o cálculo
elemento a elemento e o Pillow gera o PNG colorido, que volta em base64
e o Leaflet sobrepõe ao mapa.

A fórmula é \textbf{NDVI = (NIR $-$ Red) / (NIR + Red)}, clipada em
[$-$1, 1]. O gradiente vai de marrom (solo) para amarelo (estresse) até
verde (vegetação saudável).\,\textquotedblright

\subsection*{Slide 10 — Demonstração \hfill \textit{João V. Castello · 3--4 min}}

\textquotedblleft\,Agora vamos ver isso funcionando. Eu vou conduzir a
demo enquanto a equipe complementa.\,\textquotedblright

\textbf{Roteiro da demo:}
\begin{enumerate}[label=\textbf{0\arabic*}]
  \item Carregar o talhão \texttt{crop\_field\_1} em Picos-PI.
  \item Visualizar o \textbf{NDVI} — overlay no mapa, estatísticas e histograma.
  \item Clicar em um pixel — popup com nível de estresse na coordenada.
  \item Trocar para \textbf{EVI} e comparar lado a lado.
  \item Ativar o \textbf{terreno 3D} e rotacionar (Three.js).
  \item Acionar a \textbf{classificação em zonas} (K-Means) — hectares por classe.
\end{enumerate}

\textbf{Fala de fechamento da demo:} \textquotedblright\,Tudo isso roda
localmente, sem login, sem nuvem, e em alguns segundos por imagem.
Vou passar para o Lauan falar sobre testes.\,\textquotedblright

\subsection*{Slide 11 — Estratégia de testes \hfill \textit{Lauan Matheus · 75 s}}

\textquotedblleft\,Sobre testes, montamos uma \textbf{pirâmide em três
níveis}.

\textbf{Nível 1 — Unitário.} pytest direto no
\texttt{SentinelProcessor}: NDVI, EVI, estatísticas, casos de borda
como NaN, divisão por zero e máscara de água. Esse nível roda em
milissegundos.

\textbf{Nível 2 — Integração.} FastAPI TestClient valida os endpoints
HTTP, os status codes e a validação Pydantic do contrato.

\textbf{Nível 3 — Manual.} Testes end-to-end com um produto Sentinel-2
real, executados durante a demonstração que vocês acabaram de ver.

Tudo isso é viabilizado por \textbf{fixtures sintéticas em NumPy}, que
substituem um produto Sentinel real de 1\,GB. A suíte completa roda em
\textbf{menos de 1 segundo}.\,\textquotedblright

\subsection*{Slide 12 — Resultados \hfill \textit{Lauan Matheus · 75 s}}

\textquotedblleft\,Os números: \textbf{28 de 28 testes passando} em
\textbf{0,79 segundos}, com \textbf{zero \textit{flaky}}. A cobertura
está dividida por camada: \textbf{55\%} no
\texttt{sentinel\_processor}, que é onde a lógica de cálculo mora;
e \textbf{24\%} global, porque os endpoints exigem um produto
Sentinel real para serem exercitados.

E isso é \textbf{honesto}: declaramos esse gap como
\textbf{dívida técnica DT-03} no documento de entrega.
Rastreabilidade: os requisitos \textbf{RF-04, 05, 06, 10, 11 e 20}
estão diretamente cobertos por casos de teste vinculados.\,\textquotedblright

\subsection*{Slide 13 — Limitações + próximos passos \hfill \textit{Sammuel Moura · 90 s}}

\textquotedblleft\,Encerro o bloco técnico com honestidade sobre
\textbf{limitações conscientes}:
\begin{itemize}[noitemsep]
  \item Talhões desenhados \textbf{não persistem} — o estado é em memória.
  \item A classificação supervisionada está \textbf{simulada} (K-Means + heurística).
  \item As imagens Sentinel-2 são \textbf{baixadas manualmente} do Copernicus Hub.
  \item Aplicação \textbf{single-user} em localhost, sem autenticação.
  \item A cobertura de I/O depende de produto \texttt{.SAFE} real.
\end{itemize}

E \textbf{sete próximos passos} concretos:
\textbf{(01)} integração com a API do Copernicus para baixar imagem
sob demanda; \textbf{(02)} persistência em SQLite ou PostGIS;
\textbf{(03)} Random Forest para classificação real;
\textbf{(04)} exportação em GeoTIFF, Shapefile e CSV;
\textbf{(05)} comparação temporal multi-data;
\textbf{(06)} hospedar a aplicação para multi-usuário; e
\textbf{(07)} CI automatizado com pytest + GitHub Actions.

Vou passar para o João Leonardi encerrar.\,\textquotedblright

\subsection*{Slide 14 — Equipe + Perguntas \hfill \textit{João Leonardi · 60 s}}

\textquotedblleft\,Para fechar, a divisão de trabalho da equipe:
eu fiquei no \textbf{backend FastAPI e integração Git};
João Castello no \textbf{frontend e na condução desta demo};
José Melquíades na \textbf{modelagem de dados e parsing KML};
Lauan Matheus na \textbf{suíte de testes};
Lucas Benevinuto na \textbf{apresentação};
Luis Gustavo na \textbf{documentação técnica e ADRs};
Sammuel Moura no \textbf{manual do usuário e glossário};
e Vinicius Henrique em \textbf{arquitetura, DevOps e roteiro de demo}.

\textbf{Obrigado pela atenção.} Estamos abertos a perguntas.\,\textquotedblright

\newpage

\section*{Dicas de execução}

\begin{itemize}
  \item \textbf{Transições.} Cada apresentador encerra com uma frase que
        \textbf{passa a bola} explicitamente (\textit{``vou passar para
        o X falar sobre Y''}). Evita pausas mortas e cola o pitch.
  \item \textbf{Demo (slide 10).} O condutor deve ter o produto
        \textbf{pré-aberto e com o KML já carregado}. O risco maior é
        gastar tempo procurando arquivo na hora.
  \item \textbf{Números que vocês precisam saber de cabeça:}
        \textbf{35 histórias} (\textbf{30 entregues}, 85{,}7\%),
        \textbf{11 ADRs}, \textbf{28/28 testes em 0{,}79\,s},
        cobertura \textbf{55\% processador / 24\% global},
        \textbf{6 épicos}, \textbf{5 \textit{Won't Have}} conscientes.
  \item \textbf{Se perguntarem por que não tem banco}: a resposta está
        em DA-02 — comunicação por JSON+base64, sem persistência, para
        manter o sistema reprodutível e single-user em localhost.
  \item \textbf{Se perguntarem sobre cobertura baixa de 24\%}: é um
        número honesto — declarado como DT-03 — e é consequência direta
        dos endpoints precisarem de um produto Sentinel \texttt{.SAFE}
        real, que tem 1\,GB. O que importa é os \textbf{55\% no
        processador}, que é onde a lógica de cálculo vive.
  \item \textbf{Ritmo.} Mantenham cerca de \textbf{1 minuto por slide}
        (exceto a demo). Se ficarem atrasados, é melhor abreviar o
        slide 7 (Stack) — é o mais denso e menos crítico para a banca.
\end{itemize}

\end{document}
